\section{Metodología}

El trabajo se desarrollará en seis fases, en correspondencia con los objetivos específicos.

\noindent\textbf{F1. Caracterización térmica y mecánica.} Banco de ensayo con termopares tipo K acondicionados por MAX31855 y termómetro infrarrojo, para registrar temperatura contra tiempo en ambas caras y a distancias laterales conocidas de la línea, en los tres espesores. El \emph{springback} se medirá con goniómetro digital sobre probetas plegadas a \SI{90}{\degree} bajo distintas combinaciones de temperatura de consigna y tiempo de sostenimiento, construyendo la superficie de respuesta que alimentará la compensación del controlador. Según disponibilidad se verificará $T_g$ por calorimetría diferencial de barrido.

\noindent\textbf{F2. Subsistema térmico.} Modelo de conducción transitoria unidimensional a través del espesor, con fronteras de radiación y convección, para estimar potencia lineal y tiempo de calentamiento. Sobre esa base se selecciona el calefactor —se parte de resistencia tubular en canal reflector de aluminio, comparada contra nicromo tensado y cinta calefactora por uniformidad, inercia térmica, costo y seguridad— y se dimensionan aislamiento y sensor. La uniformidad longitudinal se verificará con termopares distribuidos e imagen térmica.

\noindent\textbf{F3. Diseño mecánico.} Modelado CAD de estructura, sujeción de la lámina, tope de posicionamiento longitudinal motorizado y ala de plegado con motor paso a paso y reductor, dimensionado a partir del par requerido para plegar en estado gomoso estimado en F1, con factor de seguridad. Fabricación en el taller de la Universidad con perfilería de aluminio estructural.

\noindent\textbf{F4. Control e interfaz.} Identificación del conjunto calefactor--lámina como sistema de primer orden con retardo mediante ensayo de escalón, y sintonización de un PID discreto en Arduino con relé de estado sólido y protección por sobretemperatura. El plegado se programará como máquina de estados con compensación de \emph{springback} y enfriamiento bajo restricción. La interfaz táctil permitirá fijar espesor, ángulo y posición, y mostrará el estado térmico.

\noindent\textbf{F5. Integración y construcción.} Ensamble mecánico, cableado de potencia y señal en gabinete, paro de emergencia, protecciones eléctricas y guardas térmicas, y puesta en marcha con pruebas incrementales.

\noindent\textbf{F6. Validación y entrega.} Al menos \num{30} repeticiones por combinación de espesor y ángulo (\num{45}, \num{90} y \SI{120}{\degree}), midiendo ángulo con goniómetro digital y posición con calibrador; se reportarán media, desviación estándar e incertidumbre expandida y se evaluará la calidad óptica contra patrón. Se elaborarán el manual de operación y mantenimiento y la documentación de entrega.
